\section{A class-correlated representation error, and what it does locally}
\label{sec:mechanism}

Nothing in \S\ref{sec:model} refers to an agent's class. We now let a fraction
$\fd$ of agents make one specific mistake, and show that its local consequence
can be worked out exactly.

\paragraph{The error is in the representation of an interaction, not of an issue.}
A receiver scoring an incoming message forms a representation of the
\emph{interaction}: the issue, the opinion asserted, and --- for a discriminating
agent --- one further coordinate that depends only on who sent it. Write
$z_r\in\{0,1\}$ for the discriminator indicator, drawn independently of class
with $\mathbb{P}(z_r{=}1)=\fd$, and
%
\begin{equation}
  \varphi_r(\hat{x},\sigma_e,\kappa_e) \;=\;
  \Bigl[\; \sigma_e\hat{x}/\gamma_C \;\;;\;\; z_r\,s(\kappa_r,\kappa_e) \;\Bigr]
  \in\mathbb{R}^{K+1},
  \qquad
  W_r \;=\; \bigl[\,\hat{w}_r \;;\; d\,\bigr],
  \label{eq:augmented}
\end{equation}
%
with $s(\kappa_r,\kappa_e)=\kappa_r\kappa_e$. The compatibility score
$W_r\cdot\varphi_r$ is then exactly the agreement field of \eqref{eq:fields}
shifted by a class-dependent constant,
%
\begin{equation}
  \Hf \;=\; \hw + \Dfield,
  \qquad
  \Dfield \;=\; z_r\,d\,\kappa_r\kappa_e ,
  \label{eq:disc-field}
\end{equation}
%
and the receiver evaluates \eqref{eq:evidence}--\eqref{eq:modulation} at $\Hf$ in
place of $\hw$. Two features of \eqref{eq:augmented} are load-bearing and are
easy to get wrong. Attaching the extra coordinate to the \emph{issue} rather
than the interaction would multiply the shift by $\sigma_e$ and divide it by
$\gamma_C$, so it would flip with the emitter's stated opinion and decay as the
agent grew confident --- neither of which is intended. And the coordinate is
deterministic, carrying no posterior variance, which is what leaves $\gamma_C$
untouched; a learned coordinate of variance $v_a$ would perturb the
normalisation at order $v_a/\gamma_C^{2}$. Appendix~\ref{app:representation}
gives the assumptions in full, together with the learned-prejudice variant.

\paragraph{The feature is task-irrelevant, and the population makes it relevant.}
The receiver models the label as $\sigma_T=\mathrm{sign}(w_T\cdot\hat{x})$ for an
unknown teacher $w_T$, so the class label is irrelevant in the exact sense
$I(\sigma_T;\kappa_e \mid \hat{x}, w_T)=0$: conditioned on the issue and the
teacher, knowing the source's class says nothing about the correct answer. In
the population, weights and labels are drawn independently, so this holds at
initialisation by construction. It does not continue to hold. Once opinions
align with class, the label \emph{does} predict what a source will say, and the
spurious feature becomes self-fulfillingly valid --- an agent audited against the
realised data distribution would find no bias to remove. We return to this in
\S\ref{sec:conclusion}.

\paragraph{Sign convention, stated once.}
$d>0$ denotes \textbf{tolerance towards the in-group and intolerance towards the
out-group}. This is forced by the algorithm rather than chosen: by
\eqref{eq:modulation} the factor $1-2\Phi(\Hf)$ is negative when $\Hf>0$ and
$\mu$ is \emph{distrust}, so adding $+d$ makes a receiver read its counterpart as
more agreeable than it is and grow more trusting. Adding $-d$ manufactures
disagreement and breeds distrust. Earlier presentations of this model tabulate
the opposite assignment while placing the discriminatory regime at $d>0$ in
their text and figures; the two are inconsistent, and taken literally they
exchange the halves of the diagram. Appendix~\ref{app:representation} works the
discrepancy through and shows both readings.

\subsection{The local flow: an exact symmetry, exactly broken}
\label{sec:flow}

Held at a fixed issue and emitter opinion, a receiver's two fields move along
$(\Fw,\Fmu)$ evaluated at $(\Hf,\hmu)$. Appendix~\ref{app:derivation} gives the
exact increments, which carry additional terms in $\FC$ and $\FV$; those terms
share the symmetry below and are cubic near the fixed point, so they change
neither the symmetry nor the linearisation. We therefore state the result for
the leading flow $X_D(\hw,\hmu)=\bigl(\Fw(\Hf,\hmu),\,\Fmu(\Hf,\hmu)\bigr)$.

\begin{proposition}[Rigid translation, exact symmetry, explicit separatrix]
\label{prop:sym}
Let $Z$, $\Fw$, $\Fmu$ be as in \eqref{eq:evidence}--\eqref{eq:modulation}. Then
\begin{enumerate}\itemsep1pt
  \item[(i)] \textbf{Translation.} $X_D(\hw,\hmu)=X_0(\hw+D,\hmu)$ for every $D$.
    The portrait at field $D$ is the unbiased portrait translated rigidly by
    $-D$ along the agreement axis.
  \item[(ii)] \textbf{Nullclines.} $\Fw$ vanishes exactly on $\{\hmu=0\}$ and
    $\Fmu$ exactly on $\{\Hf=0\}$, so $(-D,0)$ is the unique fixed point.
  \item[(iii)] \textbf{Symmetry at $D=0$.} $X_0$ is equivariant under the point
    reflection $u\mapsto-u$ and under exchanging the two sectors.
  \item[(iv)] \textbf{Separatrix.} The line $\hmu=\hw+D$ is globally invariant,
    and $(-D,0)$ is a hyperbolic saddle with eigenvalues $\mp2/\pi$. Below the
    line the flow runs to the consonant state \emph{agree and trust}; above it,
    to \emph{disagree and distrust}.
\end{enumerate}
\end{proposition}

The proof is three lines of $\Phi(-a)=1-\Phi(a)$ and is given in
Appendix~\ref{app:derivation}; every claim is also verified numerically in the
accompanying code. Part (i) is stronger than a perturbative statement: there is
no expansion in $d$ anywhere. Part (iii) is the symmetry that the discrimination
field breaks, and part (iv) says how --- the basin of \emph{agree and trust} is
$\{\hmu<\hw+D\}$, which grows strictly with $D$. A discriminating receiver
therefore has an enlarged trust basin towards its own class and an enlarged
distrust basin towards the other, and every such agent is biased towards the
same alignment of distrust with class.

\subsection{A prediction with no free parameters}
\label{sec:prediction}

Proposition~\ref{prop:sym}(iv) says a channel is captured by the trust attractor
exactly when its initial dissonance gap $\psi=\hmu-\hw$ falls below $D$. Writing
$\Psi$ for the distribution of $\psi$ at initialisation, the expected asymptotic
trust on a channel carrying field $D$ is $\Gresp(D)=2\Psi(D)-1$. Under
\eqref{eq:disc-field} with balanced classes, a non-discriminating receiver
contributes nothing on average because $\Gresp(0)=0$, so the trust--class
correlation of \S\ref{sec:setup} obeys
%
\begin{equation}
  \CcT \;=\; \fd\,\Gresp(d).
  \label{eq:prediction}
\end{equation}
%
At the model's own initialisation $\Psi$ is available in closed form: $C=I$ and
unit issues give $\gamma_C=\sqrt2$ and $\hw\sim\mathcal{N}(0,\tfrac12)$, while
$V=1$ and $\mu\sim\mathcal{U}(-1,1)$ give $\hmu$ uniform on
$\pm1/\sqrt2$, so $\psi$ is the convolution of a uniform and a Gaussian of equal
scale $1/\sqrt2$ and $\Psi$ integrates elementarily
(Appendix~\ref{app:derivation}). \emph{Equation \eqref{eq:prediction} therefore
has no fitted quantity of any kind.} It predicts that $\CcT$ is odd in $d$; that
dividing the measured map by $\fd$ collapses it onto the single curve $\Gresp$;
that $|\CcT|\le\fd$; and that the initial susceptibility is
$\partial_d\CcT|_{0}=\fd\,\Gresp'(0)$ with $\Gresp'(0)=0.9655$.

Equation \eqref{eq:prediction} is a \emph{local} theory: it treats each channel
as deciding on its own. \S\ref{sec:regimes} shows the measured correlation
exceeds it, and we report the residual $\Delta=\CcT-\fd\Gresp(d)$ as the
population's own contribution --- the amplification that comes from
non-discriminating agents acquiring class-aligned trust second-hand.

\subsection{Why distrust follows class before opinion does}
\label{sec:ordering}

The field enters the two sectors asymmetrically, and the asymmetry is structural
rather than quantitative. By \eqref{eq:modulation},
$\mathrm{sign}(\Fmu)=-\mathrm{sign}(\Hf)$, so $D$ can \emph{flip the sign} of the
trust update whenever $|\hw|<|D|$. But $\mathrm{sign}(\Fw)$ depends on $\hmu$
alone: a receiver moves towards a trusted source's opinion and away from a
distrusted one's whatever $D$ is, and $D$ only rescales the step. Class
structure therefore reaches the opinion sector only \emph{through} the trust
sector. The direct representational drive on the opinion--class correlation is
moreover of the opposite sign and roughly twenty-five times weaker --- the
relevant coefficients are $\mathbb{E}[\FC|w_e\!\cdot\!\hat{x}|]=-0.0247$ against
$\partial_{\Hf}\Fmu|_0=-2/\pi$ (Appendix~\ref{app:derivation}).

\begin{corollary}
\label{cor:ordering}
$\mathrm{sign}(\CcT)=\mathrm{sign}(d)$, with $|\CcT|\ge|\CcO|$ throughout;
$\CcO$ is small and of the \emph{opposite} sign at small $\fd|d|$, and becomes
positive only once $\fd\Gresp(d)$ is large enough for the trust field to
reorganise opinions.
\end{corollary}

The small-$|d|$ negative lobe in $\CcO$ is a falsifiable prediction that no
prior description of this model makes; \S\ref{sec:regimes} tests it.

\paragraph{There is no critical field strength for $\CcT$, and there cannot be.}
It is tempting to read the sharp sign change of $\CcT$ at $d=0$ as a transition
and to look for a critical curve $d_c(\fd)$. That search cannot succeed.
$\Fmu$ vanishes \emph{identically} on the line $\hw=0$, so
$\partial_{\hmu}\Fmu=0$ there: the class-trust contrast has no linear restoring
force at the class-symmetric state. The contrast is a marginally damped, driven
response rather than an instability, and by the class-relabelling symmetry
$\CcT$ is an odd function of $d$ with $\CcT\equiv0$ at $d=0$ for every $N$. The
sign change is a \emph{forced zero-crossing of an odd response}, not a
bifurcation. This also explains, structurally, why an earlier semi-analytic
closure of this model obtained a largest eigenvalue that was identically zero
across its whole parameter grid: the vanishing is a property of the rule, not a
shortcoming of that closure. Where a genuine threshold may live is the onset of
\emph{ideological} alignment, and \S\ref{sec:regimes} looks for it there.
